\documentclass[11pt,a4paper]{article}
\usepackage[margin=26mm,headheight=15pt]{geometry}
\usepackage[T1]{fontenc}
\usepackage[utf8]{inputenc}
\usepackage{libertinus}
\usepackage{amsmath,amssymb,amsthm,mathtools}
\usepackage{microtype,booktabs,array,longtable}
\usepackage{xcolor,enumitem,fancyhdr,titlesec,xurl}
\usepackage{aliascnt}
\usepackage{hyperref}
\usepackage[nameinlink,noabbrev]{cleveref}
\definecolor{ink}{RGB}{25,53,75}
\definecolor{muted}{RGB}{85,90,96}
\hypersetup{colorlinks=true,linkcolor=ink,citecolor=ink,urlcolor=ink,
 pdftitle={Critical Cusps in Digital-Product Pressure},
 pdfauthor={Research draft prepared with ChatGPT for Vladimir Reshetnikov},
 pdfsubject={Critical and subcritical Thue--Morse pressure; Jordan collision and finite-size scaling}}
\titleformat{\section}{\Large\bfseries\color{ink}}{\thesection}{.7em}{}
\titleformat{\subsection}{\large\bfseries}{\thesubsection}{.7em}{}
\pagestyle{fancy}\fancyhf{}
\fancyhead[L]{\small Critical cusps in digital-product pressure}
\fancyhead[R]{\small ProveIt research}
\fancyfoot[C]{\thepage}
\setlength{\emergencystretch}{3em}
\setlength{\parskip}{.2em}
\setlist{itemsep=.25em,topsep=.5em}
\allowdisplaybreaks
\numberwithin{equation}{section}
\newtheorem{theorem}{Theorem}[section]
\newaliascnt{lemma}{theorem}\newtheorem{lemma}[lemma]{Lemma}\aliascntresetthe{lemma}
\newaliascnt{proposition}{theorem}\newtheorem{proposition}[proposition]{Proposition}\aliascntresetthe{proposition}
\newaliascnt{corollary}{theorem}\newtheorem{corollary}[corollary]{Corollary}\aliascntresetthe{corollary}
\theoremstyle{definition}
\newaliascnt{definition}{theorem}\newtheorem{definition}[definition]{Definition}\aliascntresetthe{definition}
\newtheorem{question}{Research question}
\theoremstyle{remark}
\newaliascnt{remark}{theorem}\newtheorem{remark}[remark]{Remark}\aliascntresetthe{remark}
% ed. (2026-09-29): unnumbered environment for editorial notes.
\newtheorem*{ednote}{Editorial note (ProveIt, 2026-09-29)}
% ed. (2026-09-30): a second unnumbered environment for the notes of 2026-09-30.
\newtheorem*{ednotelater}{Editorial note (ProveIt, 2026-09-30)}
\newcommand{\TT}{\mathbb T}
\newcommand{\RR}{\mathbb R}
\newcommand{\CC}{\mathbb C}
\newcommand{\NN}{\mathbb N}
\newcommand{\ZZ}{\mathbb Z}
\newcommand{\dd}{\,\mathrm d}
\newcommand{\Lop}{\mathcal L}
\newcommand{\Id}{\mathrm I}
\newcommand{\Proj}{\Pi}
\newcommand{\PV}{\operatorname{PV}}
\newcommand{\spec}{\operatorname{spec}}
\newcommand{\dist}{\operatorname{dist}}
\newcommand{\norm}[1]{\left\lVert#1\right\rVert}
\newcommand{\abs}[1]{\left|#1\right|}
\newcommand{\code}[1]{\texttt{\detokenize{#1}}}
\newcommand{\Holder}{C^\alpha(\TT)}
\newcommand{\repo}{\url{https://github.com/VladimirReshetnikov/ProveIt}}

\begin{document}
% ed. (2026-09-29): no page anchor on the title page (the build reported a
% duplicate destination page.1).
\hypersetup{pageanchor=false}
\begin{titlepage}
\thispagestyle{empty}
\setlength{\parindent}{0pt}
{\small\sffamily\color{ink} PROVEIT\quad /\quad ANALYSIS AND DYNAMICAL SYSTEMS}\par
\vspace{17mm}
{\Huge\bfseries\color{ink} Critical Cusps in\\[3mm] Digital-Product Pressure\par}
\vspace{8mm}
{\Large A Jordan collision, square-root phase scaling,\\[2mm]
and explicit subcritical response\par}
\vspace{11mm}
{\large Research draft prepared for Vladimir Reshetnikov\par}
{Prepared with ChatGPT\quad\textbullet\quad 29 September 2026\par}
\vspace{8mm}\hrule\vspace{4mm}
\begin{abstract}
For digital masks in every integer base $b\ge2$, we determine the leading
phase singularity at the critical absolute-moment exponent $s=1$:
\[
 P_{b,1}(c)=\sqrt{2(b-1)\log b}\,|c|^{1/2}
            +O_{b,\eta}(|c|^{1-\eta}),\qquad 0<\eta<\tfrac12.
\]
For the binary potential $\psi_c(x)=\log\cos^2\pi(x-c)$ this is the
previously excluded pressure order $q=1/2$, since $s=2q$.
The square root comes from a genuine size-two Jordan block of the atomic
transfer operator. We construct the generalized eigenfunction explicitly
with the digamma function and prove an exponential remainder on every
$C^\alpha$, $0<\alpha<1$, by weighted quadrature rather than an assumed
Perron theorem for a singular potential.

The same reduction gives a uniform two-parameter law when
$s-1=u\sqrt{|c|}$ and an operator-norm finite-size crossover on the scale
$n\sqrt{|c|}$. Separately, for every fixed $1/2<s<1$, we prove the exact
subcritical cusp
\[
 P_{b,s}(c)=(1-s)\log b+
 (b-1)\pi s\tan(\pi s/2)|c|+o(|c|).
\]
Its coefficient follows from a universal singular-translation integral.
Together these results distinguish the critical square root, the
subcritical linear cusp, and the repository's supercritical fractional
response. Full proofs, numerical diagnostics, a proof-dependency audit,
and a further-research agenda are included.
\end{abstract}
\vfill
{\small\color{muted} Status: unrefereed research draft with ordinary mathematical proofs;
not Lean/Rocq verified. Numerical checks are not rigorous spectral enclosures.
Priority has not been independently established.}
\end{titlepage}
% ed. (2026-09-29): page anchors resume after the title page.
\hypersetup{pageanchor=true}

\begingroup\small\setlength{\parskip}{0pt}\tableofcontents\endgroup
\newpage

\section{The research target and the scope of the answer}
\label{sec:intro}

\subsection{The published question and the repository gap}
Gohlke, Kesseb\"ohmer, and Schindler study generalized Thue--Morse
measures and establish a spectral description of their pressure, including
an analytic order-two case \cite{GKS}. Gohlke, Lamprinakis, and Schmeling
prove continuity in the phase for all nonnegative orders and explicitly
ask about stronger regularity at other positive orders
\cite[Theorems~2.2--2.3 and the following discussion]{GLS}.
Their sine convention places the exceptional atomic phase at $1/2$;
our equivalent cosine convention places it at zero.

The ProveIt repository contains two immediate predecessors. The
integer-pressure report develops finite-dimensional methods at integer
orders \cite{IntegerRepo}. The fractional-pressure report proves a
sharp atomic-phase normal form for all absolute-moment exponents $s>1$
and explicitly leaves the critical and subcritical ranges untreated
\cite{FractionalRepo}. The present article addresses that excluded range.
It does not repeat the integer-order optimization theorem or the
supercritical regularity classification.

Our main target is not merely continuity at $s=1$. We determine its exact
leading singularity, its constant, its spectral mechanism, and its
finite-size transition. A second argument determines the leading phase
response for every fixed $1/2<s<1$. These are local answers at the atomic
phase, not a claim to settle phase regularity everywhere.

\subsection{What is proved here}
The principal results are the critical square-root law
(\cref{thm:critical}), the detuned spectral and operator crossover
(\cref{thm:crossover}), and the subcritical linear cusp
(\cref{thm:subcritical}). Their proofs include two structural results of
independent use: an explicit Jordan decomposition with an exponential
remainder (\cref{thm:jordan}), and a universal response formula for
translated simple zeros (\cref{lem:translation}).

The arguments rederive every repository-specific analytic ingredient
needed for these results. No unverified repository theorem is used as an
axiom. Standard facts about holomorphic functions, special functions, and
Riesz projections are identified where used; the particular perturbation
lemma needed here is proved as well.

\begin{center}
\begin{tabular}{@{}p{.25\textwidth}p{.68\textwidth}@{}}
\toprule
\textbf{Boundary} & \textbf{Precise status}\\
\midrule
Proved in this draft & Critical and detuned asymptotics in all bases;
operator finite-size scaling; subcritical response for $1/2<s<1$.\\[1.5mm]
Prior work & Atomic pressure at zero phase; positive-order continuity;
order-two analyticity; the repository's $s>1$ phase analysis.\\[1.5mm]
Not established & The phase asymptotic for $0<s\le1/2$; a global phase
regularity classification; full two-point endpoint H\"older regularity;
worldwide priority; formal verification.\\
\bottomrule
\end{tabular}
\end{center}

\subsection{Provenance and novelty limits}
The repository was inspected on 29 September 2026. The pinned head used
for provenance is
\begin{center}\small
\nolinkurl{06bcc37a82383ce99cebfec18f958487a5b3c6d4}.
\end{center}
The fractional predecessor's inspected source has Git blob identifier
\nolinkurl{83498d9ee7caa6f44b3a26b726cbb00c4b5a1a99}.
The contemporary papers \cite{GKS,GLS}, the predecessor manuscripts,
and targeted searches did not reveal a matching critical or subcritical
formula. This is a bounded provenance check, not a guarantee that these
results have never appeared elsewhere. The mathematical claim is the
proved statement below; a claim of an independently validated
``breakthrough'' would require additional scrutiny.

\section{Definitions and main theorems}
\label{sec:main}

\subsection{Masks, transfer operators, and pressure}
Fix an integer $b\ge2$ and let $T_bx=bx\pmod1$ on $\TT=\RR/\ZZ$.
Set
\begin{equation}
 d_b(x)=\frac1b\sum_{r=0}^{b-1}e^{2\pi i r x},\qquad
 m_b(x)=|d_b(x)|=
 \left|\frac{\sin\pi bx}{b\sin\pi x}\right|,
 \label{eq:mask}
\end{equation}
with the removable values at integers. For $s>0$ define
\begin{equation}
 (\Lop_{s,c}f)(x)=\sum_{k=0}^{b-1}
 m_b\left(\frac{x+k}{b}-c\right)^s f\left(\frac{x+k}{b}\right).
 \label{eq:transfer}
\end{equation}
The base is suppressed in $\Lop$ but not in the pressure.
There is \emph{no factor $1/b$} in this transfer operator.

Let $\mathcal D_{b,n}$ be the $b^n$ equal closed subintervals of $[0,1]$.
Writing $N=b^n$, put
\begin{align}
 F_{s,n,c}(x)&=\prod_{j=0}^{n-1}m_b(b^jx-c)^s,\\
 M_{s,n}(c)&=\int_0^1F_{s,n,c}(x)\dd x,\\
 Z_{s,n}(c)&=\sum_{I\in\mathcal D_{b,n}}\sup_{x\in I}F_{s,n,c}(x),\\
 P_{b,s}(c)&=\lim_{n\to\infty}\frac1n\log Z_{s,n}(c).
 \label{eq:pressure}
\end{align}
We prove existence of this last limit in all the neighborhoods under
consideration. A change of variables gives the exact normalization
\begin{equation}
 \int_0^1\Lop_{s,c}^{\,n}1\dd x=b^nM_{s,n}(c).
 \label{eq:moment}
\end{equation}
For the binary potential $\psi_c(x)=\log\cos^2\pi(x-c)$,
\begin{equation}
 P(q\psi_c)=P_{2,2q}(c).
 \label{eq:binary}
\end{equation}
In particular the critical exponent $s=1$ means $q=1/2$, not $q=1$.

For $0<\alpha<1$, $C^\alpha(\TT)$ has norm
$\norm f_\alpha=\norm f_\infty+[f]_\alpha$, with the circle distance
in its H\"older seminorm. Spectral assertions refer to its complexification.
We repeatedly use the constants
\begin{equation}
 L=\log b,\quad a_b=2\log b,\quad d=b-1,\quad
 \kappa_b=\sqrt{a_bd}=\sqrt{2(b-1)\log b},\quad
 H(x)=\frac{|\sin\pi x|}{\pi}.
 \label{eq:constants}
\end{equation}
Here $H$ is periodic, $H(x)\sim\dist(x,\ZZ)$, and
$\int_0^1H=2/\pi^2$.

\subsection{The critical square root}
\begin{theorem}[Critical phase law]
\label{thm:critical}
For every integer $b\ge2$ and every $0<\eta<1/2$, the pressure is defined
for all sufficiently small $|c|$ and
\begin{equation}
 P_{b,1}(c)=\kappa_b\sqrt{|c|}
             +O_{b,\eta}(|c|^{1-\eta}).
 \label{eq:critical}
\end{equation}
On $C^\eta$, the two eigenvalues near one are real and simple for small
$c\ne0$ and satisfy
\begin{equation}
 \lambda_\pm(c)=1\pm\kappa_b\sqrt{|c|}
                   +O_{b,\eta}(|c|^{1-\eta}).
 \label{eq:critical-eigenvalues}
\end{equation}
The remaining spectrum is in a fixed disk of radius less than one, and
$P_{b,1}(c)=\log\lambda_+(c)$.
\end{theorem}

\begin{corollary}[Sharp pointwise regularity]
\label{cor:holder}
The even function $c\mapsto P_{b,1}(c)$ has exact pointwise H\"older
exponent $1/2$ at zero. It is not Lipschitz there and has no finite
first derivative there. Its right and left difference quotients tend to
$+\infty$ and $-\infty$, respectively.
\end{corollary}
The word \emph{pointwise} is essential: this corollary asserts a bound
against the value at zero and its optimality. It does not by itself prove
a two-point $C^{0,1/2}$ bound throughout a neighborhood.

\subsection{Detuning and finite size}
Define, for $u\in\RR$,
\begin{equation}
 \omega_b(u)=\sqrt{\frac{u^2L^2}{4}+\kappa_b^2},\qquad
 r_\pm(u)=-\frac{uL}{2}\pm\omega_b(u).
 \label{eq:crossover-functions}
\end{equation}
Let $\delta_0(f)=f(0)$; the rank-one operator $H\otimes\delta_0$ sends
$f$ to $Hf(0)$.

\begin{theorem}[Uniform spectral and finite-size crossover]
\label{thm:crossover}
Fix $0<\alpha<1/2$. Let $t>0$, $s=1+ut$, and $c=\sigma t^2$, where
$\sigma\in\{-1,1\}$. Uniformly for $u$ in any fixed compact subset of
$\RR$, as $t\downarrow0$,
\begin{align}
 \lambda_\pm(s,c)&=1+t r_\pm(u)+o(t),\\
 \frac{P_{b,1+ut}(\sigma t^2)}{t}&\longrightarrow r_+(u).
 \label{eq:pressure-crossover}
\end{align}
More strongly, uniformly for $(u,\tau)$ in a compact subset of
$\RR\times[0,\infty)$, with $n=\lfloor\tau/t\rfloor$,
\begin{equation}
 t\Lop_{1+ut,\sigma t^2}^{\,n}\ \longrightarrow\
 \frac{a_b e^{-uL\tau/2}\sinh(\omega_b(u)\tau)}{\omega_b(u)}
 H\otimes\delta_0
 \label{eq:operator-crossover}
\end{equation}
in operator norm on $C^\alpha$. Consequently,
\begin{equation}
 t b^nM_{1+ut,n}(\sigma t^2)\ \longrightarrow\
 \frac{4\log b}{\pi^2}
 e^{-uL\tau/2}\frac{\sinh(\omega_b(u)\tau)}{\omega_b(u)}.
 \label{eq:moment-crossover}
\end{equation}
\end{theorem}
This identifies $n\asymp |c|^{-1/2}$ as the finite-size transition scale.
The uniformity is on compact $u$ and $\tau$ sets; divergent detuning or
divergent rescaled time is not included in the statement.

\subsection{The fixed subcritical range}
\begin{theorem}[Explicit subcritical response]
\label{thm:subcritical}
For every fixed integer $b\ge2$ and every fixed $1/2<s<1$,
\begin{equation}
 P_{b,s}(c)=(1-s)\log b+
 (b-1)\pi s\tan\left(\frac{\pi s}{2}\right)|c|+o(|c|).
 \label{eq:subcritical}
\end{equation}
Thus the two one-sided derivatives at the atomic phase are nonzero and
opposite. In the usual binary pressure order, the range covered is
$1/4<q<1/2$ and the coefficient is
$2\pi q\tan(\pi q)$.
\end{theorem}
The restriction $s>1/2$ is a rigorously identified boundary of the present
perturbation estimate. It is \emph{not} a claim that an additional
thermodynamic transition occurs at $s=1/2$.

\section{Connecting cylinder pressure to spectral growth}
\label{sec:comparison}
A transfer eigenvalue need not automatically equal a cylinder-supremum
pressure for a potential with zeros. We close that gap with a direct
polynomial argument, valid even for $0<s<1$.

\begin{lemma}[A uniform polynomial sampling inequality]
\label{lem:sampling}
Let $Q$ be a complex polynomial of degree at most $N-1$, $N\ge2$, and
let $\mathcal I_N$ be the $N$ equal cells of the circle. For every $s>0$
there is a constant $C_s$, independent of $N$ and $Q$, such that
\begin{equation}
 N\int_0^1|Q(e^{2\pi ix})|^s\dd x
 \le\sum_{I\in\mathcal I_N}\sup_I |Q(e^{2\pi ix})|^s
 \le C_sN\int_0^1|Q(e^{2\pi ix})|^s\dd x.
 \label{eq:sampling}
\end{equation}
The constants can be bounded uniformly when $s$ ranges over a compact
subset of $(0,\infty)$.
\end{lemma}
\begin{proof}
The lower bound follows by integrating on each cell. For the upper bound,
choose a maximizing point $\theta_I$ in every cell and set $R=1+1/N$.
Since $|Q|^s$ is subharmonic for every $s>0$, the Poisson inequality on
the disk of radius $R$ gives, using angles in radians,
\[
 |Q(e^{i\theta_I})|^s\le
 \int_0^{2\pi}P_{1/R}(\theta_I-\phi)|Q(Re^{i\phi})|^s
 \frac{\dd\phi}{2\pi}.
\]
The Poisson kernel satisfies
\[
 P_{1/R}(\theta)\le
 \frac{C N}{1+N^2\dist(\theta,2\pi\ZZ)^2}.
\]
There are only a bounded number of chosen points in each annular band
of angular width $1/N$ about a fixed $\phi$. This remains true if
neighboring cells choose their common endpoint. Summing the convergent
series $\sum_{j\ge0}(1+j^2)^{-1}$ yields
$\sum_I P_{1/R}(\theta_I-\phi)\le C'N$.

For $D=\deg Q$, consider the reversed polynomial
$Q^*(z)=z^D\overline{Q(1/\overline z)}$.
On each circle,
$|Q(Re^{i\phi})|=R^D|Q^*(R^{-1}e^{i\phi})|$.
The integral means of the subharmonic function $|Q^*|^s$ are increasing
with radius. Hence
\[
 \int|Q(Re^{i\phi})|^s\frac{\dd\phi}{2\pi}
 \le R^{Ds}\int|Q(e^{i\phi})|^s\frac{\dd\phi}{2\pi}
 \le e^s\int|Q(e^{i\phi})|^s\frac{\dd\phi}{2\pi}.
\]
Combining the two estimates proves the upper bound with $C_s=C'e^s$.
The zero polynomial is immediate.
\end{proof}

In our setting,
\begin{equation}
 Q_{n,c}(z)=b^{-n}\prod_{j=0}^{n-1}
 \left(\sum_{r=0}^{b-1}e^{-2\pi i r c}z^{r b^j}\right)
 \label{eq:digital-polynomial}
\end{equation}
has degree exactly $b^n-1$, and $F_{s,n,c}(x)=|Q_{n,c}(e^{2\pi ix})|^s$.
Consequently,
\begin{equation}
 b^n M_{s,n}(c)\le Z_{s,n}(c)\le C_s b^nM_{s,n}(c).
 \label{eq:pressure-comparison}
\end{equation}
Whenever the logarithmic growth rate of the integral in \cref{eq:moment}
exists, so does the pressure, and their rates agree.

At zero phase the product telescopes:
\begin{equation}
 F_{s,n,0}(x)=\left|\frac{\sin\pi b^n x}{b^n\sin\pi x}\right|^s.
 \label{eq:telescoping}
\end{equation}
Splitting into the $N=b^n$ cells and comparing distance from the two
endpoints gives
\begin{equation}
 NM_{s,n}(0)\asymp_s
 \begin{cases}
 N^{1-s},&0<s<1,\\
 1+\log N,&s=1,\\
 1,&s>1.
 \end{cases}
 \label{eq:atomic-growth}
\end{equation}
For example, the cells at distance comparable to $k/N$ contribute,
after multiplication by $N$, quantities comparable to $k^{-s}$;
on a fixed positive fraction of each cell the numerator is bounded away
from zero. Summing the resulting power series proves both bounds.
It follows that
\begin{equation}
 P_{b,s}(0)=\max\{(1-s)\log b,0\}.
 \label{eq:atomic-pressure}
\end{equation}
The binary version is already explicit in \cite{GKS}; this derivation
fixes the normalization and is not claimed as a new result.

\section{The atomic operator has a genuine Jordan block}
\label{sec:jordan}

\subsection{An eigenfunction and its generalized companion}
Let $\psi=\Gamma'/\Gamma$ be the digamma function, and define on $(0,1)$
\begin{equation}
 G(x)=-H(x)\bigl(\psi(x)+\psi(1-x)\bigr),\qquad G(0)=1.
 \label{eq:G}
\end{equation}
The endpoint expansion $\psi(x)=-x^{-1}-\gamma+O(x)$ shows that
$G(x)=1+2\gamma x+O(x^2)$ near zero; reflection gives the corresponding
expansion near one. Thus $G$ is a periodic Lipschitz function.

The following analytic family explains both the generalized eigenvector
and the detuning parameter. Put
\begin{equation}
 H_s=H^s,\quad \beta_s=b^{1-s},\quad
 a_s=\frac{2(1-\beta_s)}{s-1},\quad a_1=2\log b,
 \label{eq:family-constants}
\end{equation}
and, initially for $s\ne1$ close to one,
\begin{equation}
 G_s(x)=H(x)^s\left(\zeta(s,x)+\zeta(s,1-x)-\frac{2}{s-1}\right).
 \label{eq:Gs}
\end{equation}
Here $\zeta(s,x)$ is the Hurwitz zeta function. Its pole at $s=1$ has
residue one, and its constant term is $-\psi(x)$, so $G_1=G$.
The standard special-function identities used here are recorded in
\cite{DLMF}; their role is restricted to these explicit formulas.

\begin{lemma}[Exact two-dimensional atomic family]
\label{lem:atomic-family}
Fix $0<\alpha<1$. In a sufficiently small complex neighborhood of $s=1$,
$H_s$ and $G_s$ are analytic $C^\alpha$-valued functions, with
$H_s(0)=0$ and $G_s(0)=1$. Moreover,
\begin{equation}
 \Lop_{s,0}H_s=\beta_sH_s,\qquad
 \Lop_{s,0}G_s=G_s+a_sH_s.
 \label{eq:atomic-family}
\end{equation}
For real $s$ the formulas use the usual nonnegative powers.
\end{lemma}
\begin{proof}
The coboundary identity
\begin{equation}
 m_b(y)=\frac{H(T_by)}{bH(y)}
 \label{eq:coboundary}
\end{equation}
away from the zeros immediately gives the first eigenfunction equation.
For $\Re s>1$, the periodization
$S_s(x)=\sum_{j\in\ZZ}|x+j|^{-s}$ equals
$\zeta(s,x)+\zeta(s,1-x)$ on $(0,1)$. Reindexing an absolutely convergent
sum gives
\[
 \sum_{k=0}^{b-1}S_s\left(\frac{x+k}{b}\right)=b^sS_s(x).
\]
Thus $\Lop_{s,0}(H_sS_s)=H_sS_s$. Subtracting
$2H_s/(s-1)$ and using the first equation gives the second equation in
\cref{eq:atomic-family}.

For completeness, the analytic continuation is valid in the stated
Banach space, not only pointwise. Near zero, write
$S_s(x)=x^{-s}+A_s(x)$, with the pole $2/(s-1)$ removed from $A_s$.
The regularized $A_s$ is smooth in $x$ on a one-sided neighborhood and
analytic in $s$. The first part
$(H(x)/x)^s$ is smooth up to the endpoint and equals one there.
The remaining terms are $x^s$ times smooth analytic factors. For $s$
near one and $\Re s>\alpha$, these terms and all parameter derivatives
are in $C^\alpha$, uniformly on a smaller parameter disk. Factors
$(\log x)^j$ from differentiation cause no loss because
$\Re s-\alpha$ is bounded away from zero. The same argument works near
one and for $H_s$. Analytic continuation proves the identities across
$s=1$ and supplies the endpoint values.
\end{proof}

At $s=1$, \cref{eq:atomic-family} reads
\begin{equation}
 \Lop_{1,0}H=H,\qquad (\Lop_{1,0}-\Id)G=a_bH.
 \label{eq:jordan-chain}
\end{equation}
The nonzero second equation is the source of the square root.

\subsection{Dual coordinates and the exact projection}
For $s$ close enough to one, define
\begin{equation}
 \ell_s(f)=\int_0^1\frac{f(x)-f(0)G_s(x)}{H_s(x)}\dd x,
 \qquad \ell_0(f)=f(0).
 \label{eq:ell-s}
\end{equation}
The notation $\ell_0$ denotes evaluation, not the parameter value $s=0$.
We write $\ell=\ell_1$ at the critical exponent.
These are bounded functionals on $C^\alpha$: the numerator vanishes at
zero and one, and the integrand is $O(d^{\alpha-\Re s})$, where
$d=\dist(x,\ZZ)$. Choose $\Re s<1+\alpha$.
The same estimates with powers of $\log d$ prove analyticity in $s$.

By direct calculation,
\begin{equation}
 \ell_s(H_s)=1,\quad \ell_s(G_s)=0,\quad
 \ell_0(H_s)=0,\quad \ell_0(G_s)=1.
 \label{eq:dual-basis}
\end{equation}
Also $\ell_0\Lop_{s,0}=\ell_0$, since at $x=0$ all inverse-branch
weights except the branch at zero vanish. If $f(0)=0$, a change of
variables using \cref{eq:coboundary} gives
$\ell_s(\Lop_{s,0}f)=\beta_s\ell_s(f)$. Applying this to
$f-f(0)G_s$ yields
\begin{equation}
 \ell_s\Lop_{s,0}=\beta_s\ell_s+a_s\ell_0.
 \label{eq:dual-invariance}
\end{equation}
Therefore
\begin{equation}
 \Proj_sf=H_s\ell_s(f)+G_sf(0)
 \label{eq:projection}
\end{equation}
is a commuting rank-two projection. On its range, in the ordered basis
$(H_s,G_s)$, the atomic operator is exactly
\begin{equation}
 A_s=\begin{pmatrix}\beta_s&a_s\\0&1\end{pmatrix}.
 \label{eq:atomic-matrix}
\end{equation}

\subsection{Weighted quadrature and the rest of the spectrum}
It remains to show that this finite-dimensional part captures the whole
spectrum near one. A general spectral-gap theorem for a positive weight
with zeros would need hypotheses not yet checked. The explicit product
allows a direct proof instead.

\begin{lemma}[Critical weighted quadrature]
\label{lem:quadrature}
If $0<\alpha<1$, $f\in C^\alpha$, and $f(0)=0$, then
\begin{equation}
 \left\|\Lop_{1,0}^{\,n}f-
 H\int_0^1\frac{f(y)}{H(y)}\dd y\right\|_\infty
 \le C_{b,\alpha}(1+n)b^{-\alpha n}\norm f_\alpha.
 \label{eq:quadrature}
\end{equation}
\end{lemma}
\begin{proof}
Let $N=b^n$. Telescoping the weights gives, for $0<x<1$,
\begin{equation}
 \Lop_{1,0}^{\,n}f(x)=H(x)\frac1N
 \sum_{k=0}^{N-1}\frac{f((x+k)/N)}{H((x+k)/N)}.
 \label{eq:riemann-critical}
\end{equation}
Set $F=f/H$. Since $f(0)=0$, one has
$|F(y)|\le C\norm f_\alpha d(y)^{\alpha-1}$.
The first and last cells have integral $O(N^{-\alpha})$.
Their sampled terms, after multiplication by $H(x)$, have the same
uniform bound: for instance the first is at most a constant times
$N^{-\alpha}H(x)x^{\alpha-1}$, and this factor of $x$ is bounded.

On an interior cell of length $1/N$, whose distance to an endpoint is
comparable to $d\ge1/N$, the difference of two values of $F$ is bounded by
\begin{equation}
 C\norm f_\alpha
 \left(N^{-\alpha}d^{-1}+N^{-1}d^{\alpha-2}\right).
 \label{eq:cell-error}
\end{equation}
The first term uses H\"older continuity of $f$ and the second uses
$|H'|\le1$ together with $|f|\le C\norm f_\alpha d^\alpha$.
Multiplying by the cell length and summing gives
$O(\norm f_\alpha N^{-\alpha}(1+\log N))$.
Since $H$ is bounded, this proves the claimed uniform estimate.
At $x=0$ both sides of the main approximation are zero, so continuity
covers the endpoint as well.
\end{proof}

\begin{theorem}[Exact critical Jordan decomposition]
\label{thm:jordan}
On every $C^\alpha$, $0<\alpha<1$, there are commuting finite-rank
operators
\[
 \Proj=H\otimes\ell+G\otimes\ell_0,
 \qquad \mathcal N=a_bH\otimes\ell_0
\]
and an operator $R$ such that
\begin{equation}
 \Lop_{1,0}=\Proj+\mathcal N+R,\quad
 \mathcal N^2=0,\quad \Proj\mathcal N=\mathcal N\Proj=\mathcal N,
 \quad \Proj R=R\Proj=\mathcal NR=R\mathcal N=0.
 \label{eq:decomposition}
\end{equation}
For every $r$ with $b^{-\alpha}<r<1$ there is $C_r<\infty$ such that
$\norm{R^n}_{\alpha\to\alpha}\le C_rr^n$ for $n\ge1$.
Consequently, for $n\ge1$,
\begin{equation}
 \Lop_{1,0}^{\,n}f=
 H\bigl(\ell(f)+na_bf(0)\bigr)+Gf(0)+R^nf.
 \label{eq:atomic-iterate}
\end{equation}
The generalized eigenspace at one has dimension exactly two, and its
Jordan block has size exactly two.
\end{theorem}
\begin{proof}
All algebraic assertions follow from \cref{eq:jordan-chain,eq:dual-invariance}.
The invariant complement is
\[
 Y=\ker\Proj=\{f:f(0)=0,\ \ell(f)=0\}.
\]
On $Y$, \cref{lem:quadrature} gives sup-norm decay
$O((1+n)b^{-\alpha n})$. To upgrade to the full H\"older norm, first
observe from \cref{eq:jordan-chain} and the same quadrature estimate
applied to $1-G$ that
\begin{equation}
 \norm{\Lop_{1,0}^{\,m}1}_\infty\le C(1+m).
 \label{eq:linear-bound}
\end{equation}
For each fixed $m$, pairing inverse branches in the circle gives
\begin{equation}
 [\Lop_{1,0}^{\,m}f]_\alpha
 \le b^{-\alpha m}\norm{\Lop_{1,0}^{\,m}1}_\infty[f]_\alpha
       +B_m\norm f_\infty.
 \label{eq:LY}
\end{equation}
Indeed, the variation of $f$ gains the inverse-branch contraction
$b^{-\alpha m}$. The remaining term is bounded because the finitely
many product weights are Lipschitz. For points across a circle cut, the
branches are relabeled cyclically; for points separated by a fixed
positive distance, the sup norm gives the same type of bound.

Given $b^{-\alpha}<r<1$, choose $m$ large enough that
$C(1+m)b^{-\alpha m}<r^m$. Iterate \cref{eq:LY} on $Y$, inserting
\cref{eq:quadrature} into its last term. The resulting scalar convolution
of geometric sequences, with the harmless factor $1+jm$, is bounded by
a constant times $r^{jm}$. The finitely many remaining residues of $n$
modulo $m$ are absorbed into the constant. This proves the strong-norm
bound for $R^n$.

The complement therefore has spectral radius at most $b^{-\alpha}<1$.
On the range of $\Proj$, the matrix is
$\left(\begin{smallmatrix}1&a_b\\0&1\end{smallmatrix}\right)$,
with $a_b\ne0$. No additional generalized eigenvector at one can occur
in the complement. This proves the last assertion.
\end{proof}

\subsection{An exact critical normalization check}
The functional on the constant function can also be evaluated:
\begin{equation}
 \ell(1)=\int_0^1\left(
 \frac{\pi}{\sin\pi x}+\psi(x)+\psi(1-x)\right)\dd x
 =2\log\frac2\pi.
 \label{eq:ell-one}
\end{equation}
To verify the cancellation, integrate from $\varepsilon$ to
$1-\varepsilon$. The cosecant integral is
$-2\log\tan(\pi\varepsilon/2)$, while the digamma terms give
$2(\log\Gamma(1-\varepsilon)-\log\Gamma(\varepsilon))$.
The endpoint expansions give \cref{eq:ell-one}.
Integrating \cref{eq:atomic-iterate} at $f=1$ yields
\begin{equation}
 b^nM_{1,n}(0)=\frac{4\log b}{\pi^2}n+C_*+O(r^n),
 \quad
 C_*=\int_0^1G+\frac4{\pi^2}\log\frac2\pi.
 \label{eq:critical-atomic-moment}
\end{equation}
This linear growth is the finite-size signature of the Jordan block.
Its leading coefficient also agrees with the classical logarithmic
integral of the telescoped Dirichlet-type kernel.

\section{Perturbing the Jordan block}
\label{sec:perturbation}

\subsection{The norm of a phase translation}
\begin{lemma}[Phase perturbation bounds]
\label{lem:translation-norm}
For fixed $b$ and $0<\alpha<1$,
\begin{equation}
 \norm{\Lop_{1,c}-\Lop_{1,0}}_{\alpha\to\alpha}
 =O(|c|^{1-\alpha}).
 \label{eq:critical-norm}
\end{equation}
For fixed $0<s<1$ and $0<\alpha<s$,
\begin{equation}
 \norm{\Lop_{s,c}-\Lop_{s,0}}_{\alpha\to\alpha}
 =O(|c|^{s-\alpha}).
 \label{eq:subcritical-norm}
\end{equation}
If $s=1+ut$, $c=\pm t^2$, and $u$ stays bounded, then uniformly
\begin{equation}
 \norm{\Lop_{s,c}-\Lop_{s,0}}_{\alpha\to\alpha}
 =O(t^{2-2\alpha}).
 \label{eq:window-norm}
\end{equation}
\end{lemma}
\begin{proof}
The mask is Lipschitz. Its translation difference has sup norm $O(|c|)$
and bounded Lipschitz seminorm. Interpolation gives a $C^\alpha$ bound
$O(|c|^{1-\alpha})$. For $0<s<1$, the weight $m_b^s$ is $s$-H\"older;
the corresponding interpolation between the sup norm and the
$s$-H\"older seminorm gives $O(|c|^{s-\alpha})$.
Multiplication is bounded on $C^\alpha$, as is composition with each
inverse branch, so these weight estimates imply the operator estimates.
For $s\ge1$ near one, the Lipschitz constants are uniformly bounded.
For $s<1$ in the stated window,
$t^{2(s-\alpha)}=t^{2-2\alpha}\exp(2ut\log t)$, and the final
exponential factor is bounded uniformly for bounded $u$.
\end{proof}

\subsection{A second-order reduction lemma}
We need the quadratic remainder in the following standard spectral
reduction; merely a first-order error would not determine the critical
coefficient. Riesz projection methods are treated systematically in
\cite{Kato}. We give the particular calculation to fix the coordinate
convention and its error.

\begin{lemma}[Compression in a Riesz-projection chart]
\label{lem:compression}
Let a bounded operator $B_0$ on a complex Banach space have a finite-rank
commuting spectral projection $P$ for an isolated spectral cluster.
Write $E=\operatorname{ran}P$, and let $A_0=B_0|_E$ in a fixed basis.
For $B=B_0+\Delta$ with $\norm\Delta$ small, let $P_\Delta$ be the Riesz
projection defined by the same contour. The map
$U=P_\Delta|_E:E\to\operatorname{ran}P_\Delta$ is an isomorphism, and
in this chart
\begin{equation}
 U^{-1}BU=A_0+P\Delta|_E+O(\norm\Delta^2).
 \label{eq:compression}
\end{equation}
The constants are uniform in a compact norm-continuous family with
uniformly separated contours and uniformly bounded resolvents.
\end{lemma}
\begin{proof}
The resolvent identity and a Neumann series on the contour give
$P_\Delta-P=O(\norm\Delta)$. The identity $P_\Delta^2=P_\Delta$, or
its linearization, shows that
\[
 P(P_\Delta-P)P=O(\norm\Delta^2).
\]
Thus $PU=\Id_E+O(\norm\Delta^2)$ and $U$ is invertible onto the nearby
spectral subspace. If $A=U^{-1}BU$, then
\[
 (PU)A=PBU=A_0(PU)+P\Delta U.
\]
Here $PB_0=B_0P$ and $P\Delta U=P\Delta|_E+O(\norm\Delta^2)$.
Multiplication by $(PU)^{-1}$ gives \cref{eq:compression}.
Uniform bounds follow from the same contour estimates.
\end{proof}

The usual separation of the complementary spectrum also follows
directly from these estimates. A circle surrounding that spectrum and
excluding the finite cluster remains a resolvent contour for small
perturbations. Its resolvent integral bounds the powers of the
complementary part uniformly by $Cr^n$ for some $r<1$ in our critical
application.

\subsection{The matrix entry that fixes the square root}
Let $\Delta_{s,c}=\Lop_{s,c}-\Lop_{s,0}$.
Evaluation at zero gives the following exact identity:
\begin{equation}
 \ell_0(\Delta_{s,c}H_s)=
 \left|\frac{\sin\pi bc}{b\pi}\right|^s
 \sum_{k=1}^{b-1}
 \left(\frac{\sin(\pi k/b)}{|\sin\pi(k/b-c)|}\right)^s.
 \label{eq:lower-left-exact}
\end{equation}
The branch $k=0$ contributes zero because $H_s(0)=0$.
The denominators in the displayed finite sum are nonzero for small $c$.
Reflection $k\mapsto b-k$ makes the sum even in $c$, so, uniformly for
real $s$ near one,
\begin{equation}
 \ell_0(\Delta_{s,c}H_s)
 =(b-1)|c|^s\bigl(1+O(c^2)\bigr).
 \label{eq:lower-left}
\end{equation}
In particular at $s=1$ this is $d|c|+O(|c|^3)$.
It is the \emph{lower-left} entry in the ordered basis $(H,G)$.
The nonzero upper-right entry $a_b$ already exists at zero phase.

\subsection{Proof of the critical theorem}
\begin{proof}[Proof of \cref{thm:critical}]
Work on $C^\alpha$ with $\alpha=\eta\in(0,1/2)$, and write
$\varepsilon=|c|^{1-\alpha}$.
By \cref{thm:jordan}, a small circle about one encloses exactly its
two-dimensional generalized eigenspace. Apply \cref{lem:compression}
with the projection \cref{eq:projection}. Its two coordinates are
$(\ell,\ell_0)$, so \cref{eq:lower-left} gives the central matrix
\begin{equation}
 A(c)=
 \begin{pmatrix}
 1+O(\varepsilon)&a_b+O(\varepsilon)\\
 d|c|+O(\varepsilon^2+|c|^3)&1+O(\varepsilon)
 \end{pmatrix}.
 \label{eq:critical-matrix}
\end{equation}
The $O(\varepsilon^2)$ in the lower-left entry is the chart correction.
All entries are real for real $c$: the contour projection can be chosen
conjugation invariant.

Putting $z=\lambda-1$, the characteristic equation is
\begin{equation}
 z^2-T(c)z-\kappa_b^2|c|
       +O(\varepsilon^2+|c|^3)=0,
 \qquad T(c)=O(\varepsilon).
 \label{eq:critical-characteristic}
\end{equation}
Because $\varepsilon=o(\sqrt{|c|})$ and
$\varepsilon^2=o(|c|)$, the discriminant is positive for small
$c\ne0$. The quadratic formula yields two distinct real roots.
The error in their leading square roots is
$O(\varepsilon^2/\sqrt{|c|}+|c|^{5/2})$, which is
$O(\varepsilon)$ when $\alpha<1/2$. This proves
\cref{eq:critical-eigenvalues}. The complementary spectrum and its
powers remain uniformly inside a disk of radius less than one.

We must still show that the larger eigenvalue actually controls the
integral of $\Lop^n1$. Let $P_\pm(c)$ be the two simple eigenprojections.
The two-dimensional formula
\[
 P_+(c)|_{\mathrm{central}}=
 \frac{A(c)-\lambda_-(c)\Id}{\lambda_+(c)-\lambda_-(c)}
\]
and its minus counterpart, together with the chart
$U=\Id+O(\varepsilon)$, give
\begin{equation}
 \int_0^1P_\pm(c)1\dd x
 \sim \pm\frac{a_b}{\pi^2\kappa_b\sqrt{|c|}}.
 \label{eq:critical-amplitudes}
\end{equation}
Indeed the coordinate of $1$ along $G$ tends to $1$, the upper-right
matrix entry tends to $a_b$, and $\int H=2/\pi^2$.
The coefficient for the larger eigenvalue is positive, not zero.
For each fixed small $c\ne0$, it follows that
\[
 \int\Lop_{1,c}^{\,n}1
 =\left(\int P_+(c)1\right)\lambda_+(c)^n
 +\left(\int P_-(c)1\right)\lambda_-(c)^n+O(r^n).
\]
Thus its logarithmic growth rate is $\log\lambda_+(c)$.
The sampling comparison \cref{eq:pressure-comparison} identifies this
rate with the defined pressure. At $c=0$ use \cref{eq:atomic-pressure}.
Finally $\log(1+z)=z+O(z^2)$ supplies \cref{eq:critical}.
\end{proof}

\begin{proof}[Proof of \cref{cor:holder}]
Reflection $f(x)\mapsto f(-x)$ conjugates the operators at $c$ and $-c$;
the cylinder definition has the same symmetry. Thus the pressure is
even. Since $P_{b,1}(0)=0$ and $\kappa_b>0$, \cref{eq:critical} gives
$P_{b,1}(c)/\sqrt{|c|}\to\kappa_b$. This proves all assertions,
including failure of every pointwise H\"older exponent greater than
$1/2$.
\end{proof}

\section{The two-parameter collision and its finite-time law}
\label{sec:crossover}

\subsection{Why an analytic parameter basis matters}
A direct comparison of $\Lop_{1+ut,c}$ with $\Lop_{1,0}$ has an
$O(t)$ exponent perturbation. Squaring this error would be of the same
size as the decisive $t^2$ lower-left entry. The exact $s$-dependent
basis in \cref{lem:atomic-family} avoids that loss: all exponent
variation is treated exactly before the phase is perturbed.

Fix $0<\alpha<1/2$ and a small neighborhood of $s=1$ with
$\alpha<\Re s<1+\alpha$. The family $s\mapsto\Lop_{s,0}$ is analytic
on $C^\alpha$. To see this, differentiate the scalar weights
$m_b^s$ in $s$; terms $m_b^s(\log m_b)^j$ remain in $C^\alpha$ on a
smaller parameter disk because $\Re s>\alpha$.
The projection $\Proj_s$ in \cref{eq:projection} is also analytic and
commuting. Its complementary part has spectrum in a fixed disk
$|z|<r<1$ for $s$ sufficiently close to one, by the gap proved at one
and a resolvent Neumann series. Hence it is precisely the nearby
central spectral projection.

\subsection{The effective generator}
Set $s=1+ut$, $c=\sigma t^2$, with bounded $u$.
By \cref{lem:translation-norm}, the phase perturbation has norm
$O(\varepsilon_t)$, where
$\varepsilon_t=t^{2-2\alpha}=o(t)$.
The compression lemma now applies relative to $\Proj_s$, uniformly in
this window. In the ordered basis $(H_s,G_s)$ the matrix is
\begin{equation}
 A_{s,c}=A_s+
 \begin{pmatrix}
 \ell_s(\Delta_{s,c}H_s)&\ell_s(\Delta_{s,c}G_s)\\
 \ell_0(\Delta_{s,c}H_s)&\ell_0(\Delta_{s,c}G_s)
 \end{pmatrix}+O(\varepsilon_t^2).
 \label{eq:window-matrix}
\end{equation}
In addition to the norm estimate, the exact lower-left formula gives
\begin{equation}
 \ell_0(\Delta_{s,c}H_s)
 =d t^2\exp(2ut\log t)(1+O(t^4))=dt^2+o(t^2)
 \label{eq:window-lower-left}
\end{equation}
uniformly for bounded $u$. Notice that
$\varepsilon_t^2=o(t^2)$ is exactly where $\alpha<1/2$ is used.

Introduce the singular coordinate scaling
$D_t=\operatorname{diag}(t,1)$. Since
$\beta_{1+ut}=1-uLt+O(t^2)$ and $a_{1+ut}=a_b+O(t)$, equations
\cref{eq:window-matrix,eq:window-lower-left} imply
\begin{equation}
 D_tA_{s,c}D_t^{-1}=\Id+tJ_b(u)+o(t),\qquad
 J_b(u)=\begin{pmatrix}-uL&a_b\\d&0\end{pmatrix}.
 \label{eq:generator}
\end{equation}
All errors are uniform on compact $u$ sets. The characteristic
polynomial of $J_b(u)$ is
\begin{equation}
 z^2+uLz-\kappa_b^2,
 \label{eq:generator-polynomial}
\end{equation}
so its eigenvalues are $r_\pm(u)$ from
\cref{eq:crossover-functions}. They are distinct for every real $u$;
one is positive and the other negative.

\begin{proof}[Proof of the spectral assertions in \cref{thm:crossover}]
The two eigenvalues of \cref{eq:generator} are
$1+tr_\pm(u)+o(t)$, uniformly on compact $u$ sets. The cluster is real
and simple because its rescaled limiting roots are distinct and real.
The same projection-amplitude calculation as in
\cref{eq:critical-amplitudes} gives
\begin{equation}
 t\int P_+(s,c)1\ \longrightarrow\
 \frac{a_b}{\pi^2\omega_b(u)}>0.
 \label{eq:window-amplitude}
\end{equation}
The rest of the spectrum stays uniformly below one in modulus, while
$\lambda_+>1$ and $\lambda_-<1$ for small $t$ and bounded $u$.
It follows from \cref{eq:pressure-comparison} that the pressure equals
$\log\lambda_+$, and division by $t$ proves
\cref{eq:pressure-crossover}.
\end{proof}

\subsection{An operator, not only an eigenvalue, limit}
\begin{proof}[Proof of the finite-size assertions in \cref{thm:crossover}]
Let $n=\lfloor\tau/t\rfloor$. From \cref{eq:generator}, elementary
matrix product estimates give
\begin{equation}
 \bigl(D_tA_{s,c}D_t^{-1}\bigr)^n
 \longrightarrow e^{\tau J_b(u)}
 \label{eq:matrix-semigroup}
\end{equation}
uniformly for $(u,\tau)$ in compact subsets of
$\RR\times[0,\infty)$. One proof compares with
$(\Id+tJ_b(u))^n$ by a telescoping sum: the one-step error is $o(t)$,
there are $O(1/t)$ steps, and both matrix products are uniformly bounded
on a compact time interval. The usual exponential approximation then
finishes the argument.

For an input $f$, its coordinates on the perturbed central space are
$(\ell_s(f),f(0))+O(\varepsilon_t\norm f_\alpha)$.
After multiplication by $D_t$, these converge in operator norm to
$(0,f(0))$. On output, the map from the rescaled coordinates back to
$t$ times the original function converges to
$(v_1,v_2)\mapsto H v_1$. Indeed its columns are
$H_s+O(\varepsilon_t)$ and $tG_s+O(t\varepsilon_t)$.
Thus the limiting operator is the upper-right entry of
$e^{\tau J_b(u)}$ multiplied by $H\otimes\delta_0$.

Since
\[
 \left(J_b(u)+\frac{uL}{2}\Id\right)^2=\omega_b(u)^2\Id,
\]
that entry is
\[
 \bigl(e^{\tau J_b(u)}\bigr)_{12}
 =a_b e^{-uL\tau/2}\frac{\sinh(\omega_b(u)\tau)}{\omega_b(u)}.
\]
The complementary part contributes at most $Ct r^n$ in operator norm,
which tends to zero uniformly even when $n=0$ (then use its uniformly
bounded projection). This proves \cref{eq:operator-crossover}.
Applying the limit to $1$, integrating, and using
$\int H=2/\pi^2$ proves \cref{eq:moment-crossover}.
\end{proof}

At exact criticality $u=0$, the moment law reduces to
\begin{equation}
 t b^{\lfloor\tau/t\rfloor}
 M_{1,\lfloor\tau/t\rfloor}(\pm t^2)
 \longrightarrow
 \frac{4\log b}{\pi^2\kappa_b}\sinh(\kappa_b\tau).
 \label{eq:critical-sinh}
\end{equation}
For small rescaled time its right side is
$(4\log b/\pi^2)\tau+O(\tau^3)$, in agreement with the atomic linear
law \cref{eq:critical-atomic-moment}. For larger fixed rescaled time it
contains both exponentials coming from the split Jordan block. This
explains why replacing the whole finite-size behavior by only its
largest eigenvalue misses the crossover near the atomic phase.

\section{Subcritical response from translated zeros}
\label{sec:subcritical}

\subsection{The dominant atomic eigenvalue is simple below one}
For $0<s<1$, the eigenvalue $\beta_s=b^{1-s}>1$ has right eigenfunction
$H_s$ and a particularly simple dual functional:
\begin{equation}
 \nu_s(f)=\int_0^1\frac{f(x)}{H_s(x)}\dd x,
 \qquad \nu_s(H_s)=1.
 \label{eq:subcritical-dual}
\end{equation}
The integral is finite because $s<1$, even for $f(0)\ne0$.
By a change of variables, $\nu_s\Lop_{s,0}=\beta_s\nu_s$.

\begin{lemma}[Simple subcritical spectral gap]
\label{lem:subcritical-gap}
Let $0<s<1$ and $0<\alpha<s$. On $C^\alpha$,
\begin{equation}
 \beta_s^{-n}\Lop_{s,0}^{\,n}
 \longrightarrow H_s\otimes\nu_s
 \label{eq:subcritical-gap}
\end{equation}
exponentially in operator norm. In particular $\beta_s$ is a simple
isolated eigenvalue and all other spectrum has strictly smaller modulus.
\end{lemma}
\begin{proof}
For $0<x<1$, telescoping gives
\begin{equation}
 \beta_s^{-n}\Lop_{s,0}^{\,n}f(x)
 =H_s(x)\frac1N\sum_{k=0}^{N-1}
 \frac{f((x+k)/N)}{H_s((x+k)/N)},\qquad N=b^n.
 \label{eq:subcritical-riemann}
\end{equation}
The corresponding integral is $H_s\nu_s(f)$.
Splitting into the endpoint cells and the interior cells, as in
\cref{lem:quadrature}, yields the uniform estimate
\begin{equation}
 \left\|\beta_s^{-n}\Lop_{s,0}^{\,n}f-H_s\nu_s(f)\right\|_\infty
 \le C(1+n)b^{-n\min\{\alpha,1-s\}}\norm f_\alpha.
 \label{eq:subcritical-sup}
\end{equation}
Here is the difference from the critical calculation. The endpoint
singularity is now $d^{-s}$, which is integrable without subtracting
$f(0)$. Its endpoint contribution is $O(N^{s-1})$ after multiplication
by $H_s(x)$, uniformly in $x$. The interior variation of the numerator
contributes at most $CN^{-\alpha}\int_{1/N}^{1/2}y^{-s}\dd y$;
the denominator variation contributes at most
$CN^{-1}\int_{1/N}^{1/2}y^{-s-1}\dd y=O(N^{s-1})$.
At $x=0$, the normalized iterate is
$\beta_s^{-n}f(0)=N^{s-1}f(0)$, so the same bound holds there.
The extra factor $1+n$ in \cref{eq:subcritical-sup} is harmless and
allows one bound for all estimates.

The normalized iterates of $1$ are consequently uniformly bounded in
sup norm. For each fixed $m$, the analogue of \cref{eq:LY} has leading
seminorm coefficient at most $Cb^{-\alpha m}$ and a finite sup-norm
coefficient. The weights are $s$-H\"older, hence $\alpha$-H\"older;
this is sufficient for that estimate. On $\ker\nu_s$, combine the
seminorm recurrence with \cref{eq:subcritical-sup}. Choose $r<1$
larger than both $b^{-\alpha}$ and $b^{-(1-s)}$, and then choose $m$
large. The same geometric convolution argument as in
\cref{thm:jordan} proves exponential convergence in the full norm.
The limit is the rank-one commuting projection in
\cref{eq:subcritical-gap}, proving the spectral assertions.
\end{proof}

\subsection{A universal singular-translation integral}
The coefficient below is independent of the slopes of the zeros.
This fact is what turns the digital-mask calculation into a simple
multiple of the number of zeros.

\begin{lemma}[Response of finitely many simple zeros]
\label{lem:translation}
Let $m$ be a nonnegative periodic function, positive and $C^2$ away from
$N_z$ zeros. Near each zero $z_j$, assume that
$m(z_j+x)=|x|a_j(x)$ with $a_j$ positive and $C^2$.
For every $0<s<1$,
\begin{equation}
 \int_0^1\left[\left(\frac{m(x-c)}{m(x)}\right)^s-1\right]\dd x
 =N_zJ_s|c|+o(|c|),
 \label{eq:translation-integral}
\end{equation}
where
\begin{align}
 J_s&=\int_0^\infty
 u^{-s}\bigl((u+1)^s+|u-1|^s-2u^s\bigr)\dd u
 \label{eq:J-integral}\\
 &=\pi s\tan\left(\frac{\pi s}{2}\right)>0.
 \label{eq:J-value}
\end{align}
Values assigned at the finitely many zero denominators do not affect
the integral, whose singularities are integrable.
\end{lemma}
\begin{proof}
Choose small disjoint symmetric neighborhoods of the zeros, with their
closures still inside the stated local charts. On the complement,
ordinary Taylor expansion gives the first-order integral
$-sc\int (\log m)' +o(|c|)$.

In a neighborhood $[-r,r]$ of a zero, write
\[
 \left(\frac{m(x-c)}{m(x)}\right)^s
 =\frac{|x-c|^s}{|x|^s}
 \left(\frac{a(x-c)}{a(x)}\right)^s.
\]
Uniformly on this fixed neighborhood,
\[
 \left(\frac{a(x-c)}{a(x)}\right)^s
 =1-sc\frac{a'(x)}{a(x)}+O(c^2).
\]
The ratios $|x-c|^s/|x|^s$ have bounded integrals and converge to one in
$L^1([-r,r])$. Hence the local contribution is
\begin{equation}
 \int_{-r}^{r}\left(\frac{|x-c|^s}{|x|^s}-1\right)\dd x
 -sc\int_{-r}^{r}\frac{a'(x)}{a(x)}\dd x+o(|c|).
 \label{eq:local-translation}
\end{equation}
In the first integral substitute $x=|c|u$ and pair positive and negative
$u$. The result divided by $|c|$ tends to \cref{eq:J-integral}.
The paired integrand is $O(u^{-s})$ near zero and $O(u^{-2})$ at
infinity, so its integral converges throughout the stated range.
The sign of $c$ does not change this paired model integral.

All the regular first-order contributions cancel. Indeed, in a local
chart $(\log m)'=1/x+a'/a$, and the principal value of $1/x$ on a
symmetric interval is zero. Periodicity gives
\[
 \PV\int_0^1(\log m)'\dd x=0.
\]
More explicitly, after deleting symmetric intervals of radius
$\varepsilon$ around the zeros, the sum of endpoint logarithms has
terms $\log(\varepsilon a_j(-\varepsilon))-
\log(\varepsilon a_j(\varepsilon))$, each tending to zero.
Combining \cref{eq:local-translation} with the complement thus proves
\cref{eq:translation-integral} with coefficient $N_zJ_s$.

To evaluate the coefficient, substitute $u=1/v$ in
\cref{eq:J-integral}, and integrate by parts. Both endpoint terms
vanish, yielding
\begin{equation}
 \frac{J_s}{s}=
 \int_0^\infty
 \frac{(1+v)^{s-1}+\operatorname{sgn}(v-1)|v-1|^{s-1}}{v}\dd v.
 \label{eq:J-parts}
\end{equation}
The divergent pieces at zero must remain paired. In the first positive
piece use $t=v/(1+v)$, and pair it with the negative piece on $(0,1)$.
Their sum is
\[
 \int_0^1\frac{(1-t)^{-s}-(1-t)^{s-1}}{t}\dd t
 =\psi(s)-\psi(1-s)=-\pi\cot(\pi s).
\]
The remaining positive piece is
\[
 \int_1^\infty\frac{(v-1)^{s-1}}{v}\dd v
 =B(s,1-s)=\pi\csc(\pi s).
\]
The last identities are the beta and digamma reflection identities
\cite[\S5.5]{DLMF}. Therefore
$J_s=\pi s(\csc\pi s-\cot\pi s)
=\pi s\tan(\pi s/2)$, proving \cref{eq:J-value}.
\end{proof}

\subsection{Converting the integral response to pressure}
\begin{proof}[Proof of \cref{thm:subcritical}]
Fix $1/2<s<1$ and choose
\begin{equation}
 0<\alpha<s-\tfrac12.
 \label{eq:subcritical-alpha}
\end{equation}
The simple eigenvalue in \cref{lem:subcritical-gap} can be perturbed on
$C^\alpha$. The rank-one case of \cref{lem:compression}, together with
\cref{eq:subcritical-norm}, gives
\begin{equation}
 \lambda_s(c)=\beta_s+
 \nu_s\bigl((\Lop_{s,c}-\Lop_{s,0})H_s\bigr)
 +O(|c|^{2(s-\alpha)}).
 \label{eq:simple-response}
\end{equation}
The exponent in the error is greater than one by
\cref{eq:subcritical-alpha}.

Using the coboundary \cref{eq:coboundary} and changing variables over
the inverse branches, one obtains the exact response identity
\begin{equation}
 \nu_s\bigl((\Lop_{s,c}-\Lop_{s,0})H_s\bigr)
 =\beta_s\int_0^1
 \left[\left(\frac{m_b(y-c)}{m_b(y)}\right)^s-1\right]\dd y.
 \label{eq:response-exact}
\end{equation}
All singularities in this integral are integrable because $s<1$.
The mask $m_b$ has exactly $b-1$ simple zeros on the circle, at
$1/b,\ldots,(b-1)/b$. Its absolute-value local form has positive smooth
factors, so \cref{lem:translation} applies. Consequently
\begin{equation}
 \lambda_s(c)=\beta_s\bigl(1+(b-1)J_s|c|+o(|c|)\bigr).
 \label{eq:subcritical-eigenvalue}
\end{equation}
The spectral projection of the constant function has positive integral
at $c=0$, namely $\nu_s(1)\int H_s>0$, and retains this property for
small $c$. All other spectrum stays strictly smaller in modulus.
Therefore the logarithmic growth rate of $\int\Lop_{s,c}^{\,n}1$ is
$\log\lambda_s(c)$. The pressure comparison identifies this with
$P_{b,s}(c)$. Taking logarithms in
\cref{eq:subcritical-eigenvalue}, using
$\log\beta_s=(1-s)\log b$ and \cref{eq:J-value}, proves the theorem.
\end{proof}

\subsection{Why this proof stops at \texorpdfstring{$s=1/2$}{s = 1/2}}
The scalar translation integral \cref{eq:translation-integral} remains
valid for every $0<s<1$. The missing step for $s\le1/2$ is not the
integral's evaluation, nor the simple atomic spectral gap. It is the
comparison of the quadratic operator remainder with the proposed linear
term. On $C^\alpha$ the available bound is
$O(|c|^{2(s-\alpha)})$, and no positive $\alpha$ makes this $o(|c|)$
when $s\le1/2$.

This isolates a concrete research problem: replace the generic norm-square
remainder by a sharper estimate that retains the localization of the
translated zeros. The present paper does not infer a linear response
below the threshold merely from its formally correct first-order
coefficient. In particular, the formula at $q=1/4$ is not claimed.

\section{The phase diagram and what the crossover explains}
\label{sec:interpretation}

\subsection{Three different responses at the same phase}
The results can be compared without changing conventions:
\begin{center}
\begin{tabular}{@{}p{.18\textwidth}p{.47\textwidth}p{.27\textwidth}@{}}
\toprule
\textbf{Exponent} & \textbf{Leading local information} & \textbf{Source/status}\\
\midrule
$1/2<s<1$ &
$P(c)-P(0)\sim(b-1)\pi s\tan(\pi s/2)|c|$ &
\Cref{thm:subcritical}.\\[2mm]
$s=1$ &
$P(c)\sim\sqrt{2(b-1)\log b}\sqrt{|c|}$ &
\Cref{thm:critical}.\\[2mm]
$s>1$ &
$P(c)=s\log D_b(c)+2(b^s-1)\zeta(s)|c|^s+o(|c|^s)$ &
Repository predecessor \cite{FractionalRepo}.\\
\bottomrule
\end{tabular}
\end{center}
Here $D_b(c)=\sin(\pi bc)/(b\sin\pi c)$ is positive near zero.
In the third row the analytic term must not be dropped: when $s>2$ it
can dominate the displayed fractional term, and at even integer $s$
coefficient cancellations are important. That row records the
predecessor's statement, not a theorem newly reproved in this article.

The atomic spectral picture distinguishes the first two rows.
Below one, $\beta_s>1$ is a simple dominant eigenvalue with integrable
dual density $H^{-s}$. At one, that density reaches its integrability
threshold and a generalized eigenvector is required instead. The two
central eigenvalues $\beta_s$ and $1$ coalesce, while the upper-right
entry has the nonzero limit $2\log b$. Opening a phase creates a
lower-left entry of size $(b-1)|c|$. Their product produces the
quadratic equation $z^2\sim2(b-1)\log b\,|c|$.

\subsection{Matching the singular coefficients}
As $s\uparrow1$,
\begin{equation}
 (b-1)\pi s\tan(\pi s/2)\sim\frac{2(b-1)}{1-s}.
 \label{eq:subcritical-divergence}
\end{equation}
The predecessor's coefficient has the corresponding pole from above:
\begin{equation}
 2(b^s-1)\zeta(s)\sim\frac{2(b-1)}{s-1}
 \qquad(s\downarrow1).
 \label{eq:supercritical-divergence}
\end{equation}
These divergent fixed-exponent coefficients signal that those
asymptotics cannot be uniform across the critical point.
The crossover function has the consistent tails
\begin{align}
 r_+(u)&\sim\frac{2(b-1)}{u}&& (u\to+\infty),\\
 r_+(u)&=-u\log b+\frac{2(b-1)}{-u}+O(|u|^{-3})
 &&(u\to-\infty).
 \label{eq:matching}
\end{align}
This is an algebraic consistency check between the formulas. It is not
an interchange-of-limits theorem: \cref{thm:crossover} was proved only
for bounded $u$. Establishing a uniform overlap region with $u\to\infty$
would require additional remainder estimates.

\subsection{Which part is universal?}
Two constants have transparent origins. The factor $b-1$ counts the
simple zeros that feed the atomic fixed-point branch when the mask is
translated. The factor $2\log b$ is the nilpotent coupling in the
critical Jordan chain. The subcritical constant $J_s$ is a local
translation invariant independent of the slopes of those zeros.

This does not establish universality for arbitrary masks: a different
mask need not possess the coboundary \cref{eq:coboundary}, the same
critical exponent, or a two-dimensional Jordan block. What the proof
provides is a testable mechanism. If an analogous isolated Jordan chain,
a sufficiently small perturbation remainder, and a nonzero lower-left
response can be proved for another model, the same finite-dimensional
calculation determines its leading splitting.

\section{Reproducible numerical diagnostics}
\label{sec:verification}

\subsection{What the code does and does not verify}
The accompanying \code{code/verify.py} has four roles. It checks explicit
special-function identities at high precision; verifies the exact
characteristic polynomial of the $2\times2$ generator with symbolic
algebra; computes transfer-matrix collocation diagnostics; and tests the
finite-size crossover by direct matrix iteration. Its actual output is
included as \code{data/verification.json} and \code{data/run.log}.

The collocation matrix uses a uniform periodic grid and piecewise-linear
interpolation of inverse-branch evaluations. It is sparse and
nonnegative. Leading eigenvalues are computed by ARPACK with a fixed
starting vector, with residuals recorded. The main critical grid has
$65\,536$ points. Subcritical rows with $|c|=10^{-5}$ use $262\,144$
points because dividing a pressure difference by $|c|$ amplifies grid
error substantially.

\emph{These are diagnostics, not certified enclosures of the spectrum of
an infinite-dimensional operator.} A small matrix residual certifies
neither interpolation accuracy nor a limiting theorem. None of the
analytic proofs depends on the output of this program.

\subsection{Critical scaled pressure}
The values in the third column below approximate
$P_{b,1}(c)/\sqrt{|c|}$. The last column is the theorem's predicted limit.
\begin{center}
\small% Generated by code/verify.py; collocation is diagnostic only.
\begin{tabular}{@{}rrr r@{}}
\toprule
$b$ & $|c|$ & $P_{b,1}(c)/\sqrt{|c|}$ & Predicted limit\\
\midrule
2 & $10^{-2}$ & 0.90534238 & 1.17741002\\
2 & $10^{-3}$ & 1.02347259 & 1.17741002\\
2 & $10^{-4}$ & 1.10259590 & 1.17741002\\
2 & $10^{-5}$ & 1.14556562 & 1.17741002\\
3 & $10^{-2}$ & 1.58772057 & 2.09629415\\
3 & $10^{-3}$ & 1.80542336 & 2.09629415\\
3 & $10^{-4}$ & 1.95254663 & 2.09629415\\
3 & $10^{-5}$ & 2.03412370 & 2.09629415\\
5 & $10^{-2}$ & 2.66625974 & 3.58824516\\
5 & $10^{-3}$ & 3.05174974 & 3.58824516\\
5 & $10^{-4}$ & 3.31691329 & 3.58824516\\
5 & $10^{-5}$ & 3.46900368 & 3.58824516\\
\bottomrule
\end{tabular}

\end{center}
The approach to the limit is visibly slow, especially at larger base.
This is compatible with the proved error but does not measure its
implicit constant. The data are reported as computed, rather than
extrapolated values presented as certified constants.

\subsection{Subcritical linear response}
For the binary base, the corresponding diagnostic is
$(P_{2,s}(c)-(1-s)\log2)/|c|$:
\begin{center}
\small\begin{tabular}{@{}rrr r@{}}
\toprule
$s$ & $|c|$ & $(P_{2,s}(c)-(1-s)\log2)/|c|$ & Predicted limit\\
\midrule
0.60 & $10^{-3}$ & 2.4819902 & 2.5944188\\
0.60 & $10^{-4}$ & 2.5778127 & 2.5944188\\
0.60 & $10^{-5}$ & 2.5952322 & 2.5944188\\
0.75 & $10^{-3}$ & 5.2217464 & 5.6883567\\
0.75 & $10^{-4}$ & 5.6234735 & 5.6883567\\
0.75 & $10^{-5}$ & 5.6832576 & 5.6883567\\
0.90 & $10^{-3}$ & 13.2777061 & 17.8517118\\
0.90 & $10^{-4}$ & 17.0593987 & 17.8517118\\
0.90 & $10^{-5}$ & 17.7604812 & 17.8517118\\
\bottomrule
\end{tabular}

\end{center}
The $s=0.6$ row at the smallest phase slightly overshoots its limit;
this can reflect finite-phase corrections and grid error. Near $s=1$
the fixed-exponent coefficient is large, and the critical crossover
makes convergence slower. These observations are precisely why the
numerical data are not used as a substitute for the integral and
perturbation proofs.

\subsection{Identity checks and finite-size behavior}
At 70-digit working precision, the run checked
\cref{eq:atomic-family} at 120 combinations of base, exponent, and
position, including points close to the endpoints. The maximum absolute
residual was below $1.3\times10^{-67}$. The independent cutoff integral
for \cref{eq:ell-one} differed from $2\log(2/\pi)$ by about
$2.31\times10^{-15}$, consistent with its endpoint cutoff of $10^{-15}$.
The beta/digamma representation of $J_s$ was checked at $s=0.6,0.75,0.9$;
raw quadrature becomes less accurate near its integrability boundary,
and the largest observed absolute discrepancy was about
$3.18\times10^{-7}$. The exact symbolic identity checked was
$\det(z\Id-J_b(u))=z^2+uLz-a_bd$.

For $b=2$, $u=0$, and $\tau=1$, the finite-size limit in
\cref{eq:critical-sinh} is approximately $0.3504791$.
Using $c=1/n^2$, the computed values of
$n^{-1}2^nM_{1,n}(1/n^2)$ at $n=50,100,200,400$ were approximately
$0.3443066$, $0.3421737$, $0.3435866$, and $0.3458026$.
The sequence need not approach its limit monotonically. The full output
also contains detunings $u=-1,1$, pressure crossover checks at
$u=-2,-1,0,1,2$, and mesh-doubling comparisons.

A representative mesh comparison at phase $10^{-5}$ changes the
computed pressure by a few $10^{-7}$. Such a change becomes a few
hundredths after division by $|c|$. Therefore many displayed digits are
useful for reproducing the computation but should not be interpreted as
proven digits of the limiting continuum pressure.

\subsection{Rebuilding the diagnostics}
From the package root, a full run is
\begin{verbatim}
python -m pip install -r requirements.txt
OPENBLAS_NUM_THREADS=1 python code/verify.py --full
\end{verbatim}
The second line is a POSIX shell example; on Windows set the environment
variable using the shell's own syntax or omit it. The variable limits
linear-algebra threading, not mathematical precision. Omitting
\code{--full} runs a smaller diagnostic set. The script records library
versions in its JSON output and regenerates the two numerical TeX tables.

\section{A further-research agenda}
\label{sec:questions}

\begin{question}[The remaining subcritical interval]
Does \cref{eq:subcritical} remain valid for $0<s\le1/2$?
The translation integral is already known there from
\cref{lem:translation}; the task is to estimate the spectral remainder
using localization rather than its global operator norm. A useful first
goal is an $o(|c|)$ bound for the second reduced-resolvent term. A
counterexample would also be informative, since the present proof does
not predict a transition at $s=1/2$.
\end{question}
% ed. (2026-09-29): this question is answered at the atomic phase by a later
% package of the repository.
\begin{ednote}
The later repository article
\path{thue-morse/Thue_Morse_Subcritical_Pressure/} (filed 2026-09-29,
unreviewed) answers this question at the atomic phase:
\cref{eq:subcritical} holds for every $0<s<1$, with remainder
$O(|c|^{1+s-\alpha})$ for every $0<\alpha<s$, uniformly on compact
exponent intervals. Its key estimate is an $O(|c|)$ bound for the
perturbation paired with the atomic left eigenfunctional. Nonzero phases
remain open.
\end{ednote}

\begin{question}[The next critical term]
Determine the next term after $\kappa_b\sqrt{|c|}$. Is there a
$|c|\log|c|$ correction, an ordinary $|c|$ correction, or both?
This requires resolving the other central matrix entries beyond the
coarse H\"older translation bound and retaining the second-order chart
correction. An explicit coefficient would explain the slow approach in
the critical numerical table.
\end{question}
% ed. (2026-09-29): this question is answered by a later package.
\begin{ednote}
The later repository article
\path{thue-morse/Thue_Morse_Critical_Pressure_Corrections/} (filed
2026-09-29, unreviewed) answers this question: both corrections occur. For
every $0<\eta<1/2$,
\[
 P_{b,1}(c)=\kappa_b\sqrt{|c|}-(b-1)|c|\log(1/|c|)
 +\bigl[b\log b+(b-1)(\gamma-1)\bigr]|c|+O(|c|^{3/2-\eta}),
\]
with $\gamma$ Euler's constant. The logarithm arises from a finite-part
dual functional paired with the exact response of the eigenfunction $H$. A
sharp next term is not claimed.
\end{ednote}

\begin{question}[A uniform matching region]
Extend \cref{thm:crossover} from bounded $u$ to a quantitatively
controlled region $|u|\le U(t)$ with $U(t)\to\infty$.
The factor $\exp(2ut\log t)$ in
\cref{eq:window-lower-left} indicates one of the natural restrictions.
A theorem of this kind would turn \cref{eq:matching} from a consistency
check into a rigorous matched asymptotic expansion.
\end{question}

\begin{question}[Two-point endpoint regularity]
Is $P_{b,1}$ actually $C^{0,1/2}$ on a whole neighborhood of zero?
\Cref{cor:holder} only compares $P(c)$ with $P(0)$. Similarly, for fixed
$1/2<s<1$, can one prove local Lipschitz regularity and smoothness away
from the cusp, rather than only the one-sided first-order expansion?
Control of translated singularities at two nearby nonzero phases is
needed to answer these stronger questions.
\end{question}

\begin{question}[Nonatomic phases]
Classify phase regularity near nonzero rational phases and near
irrational phases. The fixed-point evaluation that exposes
\cref{eq:lower-left-exact} is special to the atomic phase. Periodic
orbits of zeros may lead to different finite-dimensional obstructions,
while recurrent nonperiodic zero orbits may require a different
functional setting.
\end{question}
% ed. (2026-09-30): a later package bears on this question
\begin{ednotelater}
\path{../Periodic_Anchors_Exact_Mixed_Moment_Phase_Diagrams/} (batch~68 of
the repository's \path{docs/incoming/}) computes exactly the pressure of
products of shifted digital masks whose phases fill complete periodic
orbits of $x\mapsto bx\bmod1$, with a common exponent on each orbit,
together with every equilibrium measure and the resulting mixed-moment
exponents, in every base. The zeros of such a product are mapped by $x\mapsto bx$ onto the selected periodic orbits. It bears on this question but does not
answer it: the phases are fixed, a single mask whose phase has period at
least two is not covered, and no regularity in the phase is proved.
\end{ednotelater}

\begin{question}[Beyond consecutive digital masks]
Find hypotheses on a trigonometric mask that imply an isolated Jordan
chain at a critical moment exponent. If its nilpotent coupling is $a$
and the leading lower-left perturbation is $d|c|^\rho$, the finite
matrix predicts a leading $\sqrt{ad}|c|^{\rho/2}$. The research task is
to derive those hypotheses and prove the required spectral separation
for actual masks, rather than infer it from the matrix alone.
\end{question}

\begin{question}[Long rescaled times and higher-order semigroups]
Obtain error bounds in \cref{eq:operator-crossover} that remain useful
for $\tau$ growing with $1/t$. Determine a next-order effective matrix
and whether it produces a uniform correction to both the pressure and
the finite-size moment. This would connect the compact-time sinh law to
the single-eigenvalue regime quantitatively.
\end{question}

\begin{question}[Eigenfunctions, eigenmeasures, and equilibrium measures]
The eigenprojection amplitudes diverge like $|c|^{-1/2}$, while their
combination has a finite rescaled operator limit. Determine normalized
weak limits and boundary layers of the individual left and right
eigenvectors, and translate them into statements about equilibrium
measures. The operator limit alone does not establish a variational
classification of such measures.
\end{question}
% ed. (2026-09-29): answered, except the left boundary layer, by a later package.
\begin{ednote}
The later article
\path{thue-morse/Thue_Morse_Critical_Pressure_Corrections/} gives the first
normalized corrections of the left and right eigenvectors, a boundary layer
of the right eigenfunction on the scale $\sqrt{|c|}$, and the weak limits of
the unique equilibrium measures: $\tfrac12\delta_0+\tfrac12$ Lebesgue
measure at $s=1$ as $c\to0$, and, along $s=1+u\sqrt{|c|}$ with $u$ in a
compact set, the mixture of $\delta_0$ and Lebesgue measure with atomic
weight $\tfrac12\bigl(1+u\log b/\sqrt{u^2\log^2b+4\kappa_b^2}\bigr)$. At
$(s,c)=(1,0)$ the equilibrium measures are exactly the segment between
$\delta_0$ and Lebesgue measure. The left boundary layer remains open.
\end{ednote}

\begin{question}[Certified spectral numerics]
Replace the collocation diagnostics by explicit operator bounds and
interval arithmetic. Useful deliverables would be rigorous enclosures
for $P_{b,s}(c)$ at nonzero phases, certified constants in the error of
\cref{eq:critical}, and mesh-independent upper and lower pressure
bounds. Non-normality near the Jordan point makes residual-only
certification especially inadequate.
\end{question}

\begin{question}[A staged Lean or Rocq formalization]
Formalize first the digital-mask coboundary, the finite-dimensional
Jordan algebra, the exact evaluation
\cref{eq:lower-left-exact}, and the matrix exponential calculation.
Then formalize the weighted quadrature estimates, bounded projections,
and the resolvent compression lemma. Only after these analytic bridges
are present should the pressure theorem be declared formally verified.
This route fits ProveIt's existing separation between executable checks
and proved semantic connections.
\end{question}

\section{Conclusion}
The atomic phase is not merely a point where a fractional exponent
formula needs a limiting interpretation. At the critical exponent the
spectral geometry changes: a simple dominant branch and an atomic
branch form a defective eigenvalue. The phase perturbation opens a
coupling in the missing matrix direction, forcing a square-root
splitting with the exact constant $\sqrt{2(b-1)\log b}$.

The argument supplies more than a scalar asymptotic. It gives an explicit
Jordan chain, a directly proved complementary gap, a two-parameter
collision law, and an operator finite-size crossover. Below the critical
exponent, a different but compatible calculation yields an explicit
linear cusp for $1/2<s<1$. The remaining interval is reduced to a
specific localized perturbation problem rather than hidden inside an
unsupported extrapolation.

\appendix
\section{Proof-dependency and claim audit}
\label{app:audit}
The logical dependencies are as follows. The sampling inequality links
integrals to cylinder pressure independently of any spectral claim.
The special-function identity constructs the critical invariant plane.
Weighted quadrature and a fixed-iterate H\"older estimate prove its
complementary gap; no general Perron theorem is used to assume that gap.
A Riesz-projection chart then reduces the small-phase problem to a
matrix with a controlled quadratic remainder. The exact fixed-point
response determines the missing entry. Only after a nonzero positive
moment amplitude is proved is the leading eigenvalue identified with
pressure.

For the subcritical theorem, the order is separate: prove a simple
atomic gap, identify its dual functional, evaluate a translated-zero
integral, bound the operator remainder by $o(|c|)$, and then pass to
pressure. The scalar first-order integral is not used as a substitute
for the remainder estimate.

\begin{center}
\begin{tabular}{@{}p{.38\textwidth}p{.54\textwidth}@{}}
\toprule
\textbf{Potential failure mode} & \textbf{How it is excluded here}\\
\midrule
Confusing moment order and pressure order & $s=2q$ is fixed in
\cref{eq:binary}; the critical order is $q=1/2$.\\[1.5mm]
Using a normalized transfer operator & \Cref{eq:transfer,eq:moment}
state the unnormalized operator and its exact integral factor.\\[1.5mm]
Assuming a gap despite zeros of the weight & The gap is proved by
weighted quadrature and a H\"older recurrence.\\[1.5mm]
Losing the decisive perturbation coefficient & The compression lemma
has an explicit norm-square error, and the lower-left entry is
computed exactly.\\[1.5mm]
Confusing spectral radius with pressure & A polynomial sampling bound
and a nonzero positive moment amplitude supply the bridge.\\[1.5mm]
Extrapolating the subcritical formula too far & The theorem is restricted
to $s>1/2$, precisely where the norm-square error is $o(|c|)$.\\[1.5mm]
Treating numerics as a proof & All infinite statements have analytic
proofs; collocation is labeled diagnostic throughout.\\[1.5mm]
Claiming full H\"older regularity from a pointwise law & The regularity
corollary explicitly states only pointwise optimality at zero.\\
\bottomrule
\end{tabular}
\end{center}

\section{Reproduction and package contents}
The package contains \code{article.tex}, its compiled
\code{article.pdf}, the verification program, the actual numerical output,
two generated numerical tables, a requirements file, a Makefile, and
separate source and validation notes. No external paper or font file is
redistributed. The document uses ordinary installed TeX packages,
including Libertinus for typesetting.

To rebuild the article without rerunning the diagnostics, execute
\begin{verbatim}
latexmk -pdf -interaction=nonstopmode -halt-on-error article.tex
\end{verbatim}
The bibliography is embedded in the TeX source; no BibTeX download is
needed. The source depends on the two included table files in
\code{data/}. A full diagnostic rerun regenerates those tables before
compilation. All new results in this package remain unrefereed and
unformalized, regardless of the formal status of other ProveIt projects.

\clearpage
\phantomsection\addcontentsline{toc}{section}{References}
\begin{thebibliography}{9}
\bibitem{GKS}
P.~Gohlke, M.~Kesseb\"ohmer, and T.~I.~Schindler,
\emph{Generalized Thue--Morse measures: spectral and fractal analysis},
arXiv:2509.22109v1, 26 September 2025.
\url{https://arxiv.org/abs/2509.22109}.
The inspected version includes the atomic pressure calculation and
Theorem~2.7 on order-two analyticity.

\bibitem{GLS}
P.~Gohlke, G.~Lamprinakis, and J.~Schmeling,
\emph{On a family of singular potentials: Parameter dependence of
thermodynamic characteristics}, arXiv:2603.19001v1, 19 March 2026.
\url{https://arxiv.org/abs/2603.19001}.
Theorems~2.2--2.3 and the following discussion motivate the regularity
question addressed locally here.

\bibitem{Repo}
V.~Reshetnikov, \emph{ProveIt}, repository snapshot
\nolinkurl{06bcc37a82383ce99cebfec18f958487a5b3c6d4},
inspected 29 September 2026.
\url{https://github.com/VladimirReshetnikov/ProveIt}.

\bibitem{IntegerRepo}
\emph{Integer Pressure and a Missing Taylor Coefficient},
research draft in ProveIt, September 2026.
Repository-relative path:
\nolinkurl{Analysis/FabiusFunction/docs/semi-formalized-research-frontiers/drafts/thue-morse/Thue_Morse_Integer_Pressure/article.tex}.
Unrefereed repository predecessor; no theorem from it is assumed in the
new critical or subcritical proofs.

\bibitem{FractionalRepo}
\emph{Fractional Cusps at the Atomic Phase: Sharp noninteger regularity
of Thue--Morse pressure and sharp relaxation for digital products in
every base}, research draft prepared with ChatGPT for V.~Reshetnikov,
29 September 2026.
Repository-relative path:
\nolinkurl{Analysis/FabiusFunction/docs/semi-formalized-research-frontiers/drafts/thue-morse/Thue_Morse_Fractional_Pressure/article.tex}.
Inspected blob:
\nolinkurl{83498d9ee7caa6f44b3a26b726cbb00c4b5a1a99}.
This predecessor explicitly excludes critical and subcritical phase
asymptotics.

\bibitem{Kato}
T.~Kato, \emph{Perturbation Theory for Linear Operators},
Classics in Mathematics, Springer, 1995, reprint of the second edition.
\url{https://doi.org/10.1007/978-3-642-66282-9}.

\bibitem{DLMF}
F.~W.~J. Olver et al., editors,
\emph{NIST Digital Library of Mathematical Functions},
\S5.5 (gamma and digamma reflection identities) and
\S25.11 (Hurwitz zeta function), accessed 29 September 2026.
\url{https://dlmf.nist.gov/5.5};
\url{https://dlmf.nist.gov/25.11}.
\end{thebibliography}
\end{document}
